\label{ch:conclusion}

\section{Summary of Contributions}

This thesis presented a controlled benchmark of five gripper variants on a
single platform (Franka FR3), measuring stationary contact force as a
function of press depth and tilt angle. The principal findings are:
\begin{enumerate}
  \item Rigid parallel-jaw grippers exhibit narrow ($\sim 1$--$2$\,mm)
        working depth ranges with steep force-depth slopes ($-19$ to
        $-26$\,N/mm).
  \item Finray-effect grippers exhibit broad ($10$--$12$\,mm) working
        ranges with a flat soft plateau ($-1.2$\,N/mm in the linear
        regime).
  \item Finray contact force is strongly tilt-sensitive at the no-contact
        boundary ($53\%$ reduction at $5^{\circ}$ tilt for the boundary
        depth of $-0.5$\,mm).
  \item Finray contact force is tilt-insensitive in the soft plateau
        ($\leq 1\%$ change at $5^{\circ}$ tilt for the safe depth of
        $+10$\,mm).
  \item Finray response is direction-dependent under tilt about the
        asymmetric ($y$) axis at boundary depths ($2.4\times$ ratio
        between $\pm 2.5^{\circ}$), but symmetric in the soft plateau.
\end{enumerate}

The depth-dependent tilt robustness identified here had not been
quantitatively characterised in prior literature, to the author's knowledge,
and constitutes the principal empirical contribution.

\section{Future Work}

\begin{description}
  \item[Cross-variant tilt sweep.] The natural extension is to repeat
        Chapter~\ref{ch:experiments} on the remaining four variants. The rigid
        parallel-jaw case is of particular interest as a hypothesised
        tilt-invariant baseline.
  \item[Tilt-sensitivity model.] An analytical model relating finger
        curvature, contact compliance, and tilt-induced contact-point
        migration would connect the present empirical findings to the
        broader continuum-mechanics literature on soft contact.
  \item[Multi-fabric study.] All trials used a single cloth specimen.
        Replicating the matrix on fabrics of varying thickness, weave,
        and surface treatment would establish the applicability range
        of the depth-tilt-robustness curves identified here.
  \item[Closed-loop control.] The boundary-depth regime, where finray
        directional asymmetry dominates, motivates a tilt-compensating
        controller using inline force feedback. Implementing such a
        controller and validating it against the open-loop boundary
        baseline would be the natural follow-on for a system-design
        thesis.
  \item[Sim2real validation.] A simulation of the finray finger in a
        compliant-contact physics engine could be calibrated against the
        empirical force-depth-tilt surface from this thesis, providing
        a digital twin for the gripper variants.
\end{description}
